\cotefile{157-2}
\scanpage{83}\mspage{27}
% L1
C'est \unsure{assez} étrange -- on va voir si c'est bien vrai !

% L2
III \uline{Construction de} $\Phi_0,\Psi_0$ :
\diagram{\begin{tikzcd}[ampersand replacement=\&,column sep=small,row sep=small]
\& \underline{\mathrm{Fibsc}}_0(S)\\
\mathrm{Cat}_{/S}\arrow[ur]\arrow[dr] \&\\
\& \underline{\mathrm{Cofibsc}}_0(S)
\end{tikzcd}}
% L3-4
On doit avoir
\[
\Phi_0(X)=\Phi r C(X)\qquad \Psi_0(X)=\Psi r C(X),
\]
% L5-8
où \deleted{\unsure{r}} \deleted{\unsure{r}} $r:\underline{\mathrm{Parf}}\,S\longrightarrow\mathrm{Cat}_{/S}$ est l'inclusion canonique. En tant que foncteurs à valeurs dans $\underline{\mathrm{Fibsc}}$ resp. $\underline{\mathrm{Cofibsc}}$, cela définit déjà $\Phi_0,\Psi_0$ par
% L9
\[
\Phi_0=\Phi r C,\quad \Psi_0=\Psi r C,
\]
% L10-14
mais il faut voir que $\Phi_0,\Psi_0$ prennent leurs valeurs \deleted{\unsure{à}}dans les sous-catégories pleines à indice $0$. \deleted{\unsure{Parf} \unsure{pour}} Cela équivaut \deleted{\unsure{donc}} au lemme suivant \deleted{(que je formule pour $\Psi_0$, pour fixer les idées)}\unsure{,|;} appliqué à $Y=rC(X)$ :

% L15-17
\uline{Lemme} Soit $Y$ parfait sur $S$ (ou seulement un Hot-fibré sur $S$). Alors $\Phi(\deleted{\unsure{X}}Y)$ est dans $\underline{\mathrm{Fibsc}}_0\,S$, $\Psi(Y)$ dans $\underline{\mathrm{Cofibsc}}_0\,S$.

% L18-19
Il faut prouver \inserted{(pour $\Psi(Y)$)} que si \deleted{\unsure{g}} $s\longrightarrow t$ est une flèche dans $S$, alors
% L20
\[
\Psi(Y)_s\unsure{\mathord{-}|\longrightarrow}\Psi(Y)_t
\]
% L21-23
est une \unsure{W|w}-équivalence. Or cette flèche s'identifie à $\deleted{\unsure{X}}Y_{/s}\longrightarrow Y_{\deleted{\unsure{\vert}}/t}$ et on sait bien que c'est \deleted{un \unsure{Ob|OK}, par définition} \inserted{une Hot-équi}%
% archive sheet number: 82
\scanpage{84}\mspage{28}%
% L1-2
valence, par définition des Hot-fibrés, donc c'est une $W$-équivalence.

% L3-4
On a donc les foncteurs
\diagram{\begin{tikzcd}[ampersand replacement=\&,column sep=large]
\deleted{\unsure{\mathrm{Cat}_{/S}\longrightarrow}}\ \underline{\mathrm{Fibsc}}_0\,S
  \arrow[r,shift left,"p"] \& \mathrm{Cat}_{/S}\arrow[l,shift left,"\Phi_0"],
\end{tikzcd}}
% L5-8
il faut voir qu'ils satisfont aux conditions équivalentes du lemme p.~11. (et ce sera pareil pour \deleted{\unsure{p}}, $q,\Psi_0$). On sait déjà que $p$ est compatible avec les localiseurs, \deleted{et $\Phi_0$}
% L9-12
\deleted{aussi car il suffit de le voir pour le foncteur \unsure{(1)} dans $\underline{\mathrm{Fibsc}}\,S$ (les localiseurs s'induisent), et alors $\Phi_0=\Phi$} Prouvons le pour $\deleted{\unsure{\Psi|\Phi}}\unsure{\Phi|\Psi}_0$. Soit
% L13
\[
 f:X\unsure{\mathord{-}|\longrightarrow}Y
\]
% L14-15
dans $\mathrm{Cat}_{/S}$ une $W$-équivalence, alors $rC(f)$ l'est aussi. Il suffit de prouver ceci :

% L16
Soit \quad $g:X'\longrightarrow Y'$
% L17-20
un morphisme dans \deleted{\unsure{Hot}}$\underline{\mathrm{Fibhot}}\,S$ (p. exemple un morphisme de $\underline{\mathrm{Parf}}\,S$), qui soit une $W$-équivalence donc des $W$-équivalences sur les fibres. Alors $\deleted{\unsure{\Psi}}\Phi(g):\deleted{\unsure{\Psi}}\Phi(X')\unsure{\mathord{-}|\longrightarrow}\deleted{\unsure{\Psi}}\Phi(Y')$ est une $W$-équivalence fibre par fibre,
% L21-22
i.e. \quad $X'_{/\deleted{\unsure{S}}s}\longrightarrow Y'_{/s}$ une $W$-équiv. pour \deleted{\unsure{tout}} $s\in S$. Ça a été vu \unsure{dans|du} Cor.~2 \unsure{p.|page}~20.
% archive sheet number: 83
\scanpage{85}\mspage{29}
% L1-2
Considérons
\[
 p\Phi_0:\mathrm{Cat}_{/S}\unsure{\mathord{-}|\longrightarrow}\mathrm{Cat}_{/S}
\]
% L3-5
ça transforme donc $X$ en $\Phi C(X)$, regardé comme objet de $\mathrm{Cat}_{/\deleted{\unsure{s}}S}$ sans plus. On a un morphisme fonctoriel
% L6
\[
 X\xrightarrow{j_X} C(X)\qquad j:\mathrm{id}_{\mathrm{Cat}_{/S}}\longrightarrow C
\]
% L7-8
d'où un morphisme fonctoriel
\[
 \Phi(X)\xrightarrow{\Phi(j_X)}\Phi C(\deleted{\unsure{x}}X)\qquad
 \Phi*j:\Phi\longrightarrow\Phi\circ C,
\]
% L9-10
et également un morphisme fonctoriel
\[
 X\xrightarrow{i_X}\Phi(X)\deleted{\unsure{\longrightarrow\Phi C(X)}}
\]
% L11-12
d'où en composant
\diagram{\begin{tikzcd}[ampersand replacement=\&]
 X\arrow[r,"i_X"]\arrow[rr,bend right=28,"\alpha_X"']\&\Phi(X)\arrow[r,"{\Phi(j_X)}"]\&\Phi C(X),
\end{tikzcd}}
% L13-14
\unsure{i.e.|ou}
\[
 \alpha:\mathrm{id}_{\mathrm{Cat}_{/S}}\longrightarrow p\Phi_0
\]
% L15-19
Pour $X$ fixé dans $\mathrm{Cat}_{/S}$, on sait que $j_X$ est dans $W$, donc aussi $\Phi(j_{\unsure{x|X}})$ (quoique non fibre par fibre), cf. cor.~1 \unsure{page|p.}~20. On sait aussi que $i_X$ est dans $W$, donc aussi \deleted{\unsure{le}} le composé $\alpha_X$.

% L20-21
Considérons
\[
 \Phi_0 p:\underline{\mathrm{Fibsc}}_0\,S\longrightarrow\underline{\mathrm{Fibsc}}_0\,S.
\]
% L22-23
On part de $X$ dans $\underline{\mathrm{Fibsc}}_0\,S$, on oublie le scindage et le fait même qu'il soit \unsure{Hot-fibrant|Cat-fibrant}
% archive sheet number: 84
\scanpage{86}\mspage{30}
% L1-4
(on a appliqué $p$), on lui applique $\Phi_0$, \deleted{\unsure{on prend}} i.e. on regarde $\Phi C X$ comme objet de $\underline{\mathrm{Fibsc}}\,S$ (et $=$ \unsure{de|dans} $\underline{\mathrm{Fibsc}}_0\,S$). On veut définir un \unsure{mor}
% L5
\[
 X\xrightarrow{\beta_X}\Phi_{\unsure{0}}p(X)=\Phi C(X)\quad\inserted{\text{dans }\text{\uuline{$\underline{\mathrm{Fibsc}}_0$}}}
\]
% L6-7
fonctoriel en $X$. Or on a comme tantôt
\[
 X\mathrel{\overset{j_X}{\unsure{\mathord{-}|\longrightarrow}}} C(X)
\]
% L8-9
d'où en appliquant $\Phi$
\[
 \Phi(X)\xrightarrow{\Phi(j_X)}\Phi C(X),
\]
% L10-15
mais \unsure{regardé|regardée} cette fois comme flèche \unsure{dans|de} $\underline{\mathrm{Fibsc}}_0\,S$ \deleted{\unsure{, d'où}} (N.B. on a bien $\Phi(X)$ dans $\underline{\mathrm{Fibsc}}_0\,S$, car $X$ \deleted{étant dans} provenant de $\underline{\mathrm{Fibsc}}_0\,S$, \deleted{\unsure{donc de}} est dans $\underline{\mathrm{Fibhot}}\,S$, et on applique la remarque ``Cor.~3'' page 21.) D'autre part, on a une flèche de $\mathrm{Cat}_{/S}$
% L16
\[
 X\xrightarrow{i_X}\Phi(X),
\]
% L17-21
est-elle une flèche dans $\underline{\mathrm{Fibsc}}_0\,S$, i.e. est-elle cartésienne ? On voit tout de suite que non, qu'on a seulement, pour une flèche $u:s\longrightarrow t$ dans $S$, un homomorphisme de commutation (pas un iso) dans
% L22-24
\diagram{\begin{tikzcd}[ampersand replacement=\&,row sep=large,column sep=large]
 X_s\arrow[r,"{u_*^{X}}"]\arrow[d,"i_s"] \& X_t\arrow[d,"i_t"{name=rightside}]\\
 X_{/s}\arrow[r,"{(\mathrm{id}_X)_{/u}}"{name=bottomside},"{\substack{\Vert\\u_*^{\Phi(X)}}}"'] \& X_{/\unsure{t|f}}
 \arrow[from=bottomside,to=rightside,bend left=12,"\alpha"]
\end{tikzcd}}
% archive sheet number: 85
\scanpage{87}\mspage{31}
% L1-4
Donc il en faut rabattre ! On peut considérer $i_X$ au mieux comme une flèche de $\underline{\mathrm{Fib}}_0\,S$, et en prenant les composés on trouve \deleted{\unsure{donc}} une flèche dans $\underline{\mathrm{Fib}}_0\,S$ :
% L5-6
\diagram{\begin{tikzcd}[ampersand replacement=\&]
 X\arrow[r,"i_X"]\arrow[rr,bend right=22,no head,"{\deleted{\unsure{B}}\beta_X\ (= \alpha_X)}"']\&\Phi(X)\arrow[r,"{\Phi(j_X)}"]\&\Phi C(X),
\end{tikzcd}}
% L7-9
i.e. (regardant \inserted{maintenant} $\Phi_0$ comme foncteur à valeurs dans $\underline{\mathrm{Fib}}_0\,S$) un hom \inserted{de} foncteurs \deleted{\unsure{en $X$}}
\[
 \deleted{\unsure{B}}\beta:\mathrm{id}_{\underline{\mathrm{Fib}}_0\,S}\longrightarrow\Phi_0 p_0,
\]
% L10-13
où $p_0:\underline{\mathrm{Fib}}_0\,S\hookrightarrow\mathrm{Cat}_{/S}$ est l'inclusion canonique. \unsure{Cet|Ce} $\beta$ est dans $W$ argument par argument, et de $=$ pour $\alpha$ page 29, interprété cette fois comme
\[
 \mathrm{id}_{\mathrm{Cat}_{/S}}\longrightarrow p_0\Phi_0.
\]
% L14
Ce qu'on a prouvé finalement est ceci :

% L15-19
\noindent\begin{tabular}{|p{0.94\linewidth}@{}}
\uline{Théorème-emballage} : Soit $\unsure{\mathcal F|F}$ une sous-catégorie \inserted{pleine} de $\mathrm{Cat}_{/S}$ qui contient une des trois sous-catégories pleines $\mathrm{Fib}_0\deleted{\unsure{(}}S\deleted{\unsure{)}}$, $\underline{\mathrm{Parf}}\,S$, $\underline{\mathrm{Cofib}}_0\,S$, et \unsure{munissant|munissons} $\mathcal F$ du localiseur $W$ ``grossier'' (induit par celui de $\mathrm{Cat}_{/S}$). Alors l'inclusion
% L20-21
\[
 \mathcal F\longrightarrow\mathrm{Cat}_{/S}
\]
induit une équivalence sur les \unsure{localiseurs|localisations}.
\end{tabular}
% vertical rule beside the theorem
% archive sheet number: 86
\scanpage{88}\mspage{32}
% L1-4
Par contre, les catégories localisées des quatre catégories de fibrés scindés
\[
 \underline{\mathrm{Fibsc}}\,S,\ \underline{\mathrm{Cofibsc}}\,S,\ \underline{\mathrm{Fibsc}}\,S,\ \underline{\mathrm{Cofibsc}}\,S
\]
échappent à \unsure{nos|ces} arguments. \uline{Et pour cause !}

% L5-10
Il est totalement \unsure{inexact}, sur $S$ connexe disons \inserted{(p. ex. $S=\Delta^1$)}, que \deleted{les \unsure{obj}} pour un type d'homotopie relatif $\deleted{\unsure{X}}\xi$ sur $S$, les types induits \deleted{\unsure{les}} $\xi_s$ $(s\in S)$ \unsure{forment} un syst. local sur $S$, à valeurs dans $\mathrm{Hot}_W$, comme ce serait le cas si le th. annoncé au début était vrai.

% L11-12
Au contraire, \uline{\unsure{toute}} flèche de $\mathrm{Hot}_W$ doit provenir d'un objet de $\mathrm{HOT}_W(\Delta^1)$.

% L13-18
Les \unsure{localisées|localisés} des 9 autres catégories des diagrammes p.~8, et de \deleted{\unsure{toutes}} celles dont il est question dans le th. qu'on vient de prouver, \inserted{doivent} correspondre ``\unsure{visiblement|évidemment}'' \unsure{à la} sous-catégorie pleine de $\mathrm{HOT}_W(S)$ formée des \unsure{HOT}-coefficients loc constants sur $S$.

% L19-22
Aussi, je me rends compte que c'était absurde de m'attendre à ce que $\underline{\mathrm{Fibsc}}\,S$ et $\underline{\mathrm{Cofibsc}}\,S$ donnent des \inserted{cat.} localisées canoniquement équivalentes. Car \inserted{les éléments de} la première
% archive sheet number: 87
\scanpage{89}\mspage{33}
% L1-3
des deux catégories donnent naissance \inserted{chacun} à \deleted{\unsure{des}} \unsure{foncteurs|foncteur} \unsure{contravariants|contravariant} sur $S$
\[
 \bar{\xi}:S^{\unsure{\mathrm o|\circ}}\longrightarrow\mathrm{Hot}_W,
\]
% L4-6
\inserted{ceux de} la seconde (qui n'est autre que $(\mathrm{Fibsc}(S^{\mathrm o}))^{\mathrm o}$) donnent naissance \inserted{chacun} à \deleted{\unsure{des}} foncteurs \uline{covariants}
\[
 \bar{\eta}:S\longrightarrow\mathrm{Hot}_W,
\]
% L7-11
ce qui n'est pas pareil ! Ce n'est ``pareil'' que si ce sont des coeff. locaux, ce qui signifie exactement que $\xi$ est dans $\underline{\mathrm{Fibsc}}_0\,S$ resp. $\eta$ dans $\deleted{\unsure{\mathrm{Fibsc}}}\underline{\mathrm{Cofibsc}}_0\,\deleted{\unsure{\xi}}S$. On s'attend donc bien \textcircled{\scriptsize 1\textsuperscript{o}} que les flèches
% L12
\[
 \underline{\mathrm{Fibsc}}_0\,\deleted{\unsure{(S)}}\unsure{X|S}\longrightarrow\underline{\mathrm{Fib}}_0\,S,\quad
 \underline{\mathrm{Cofibsc}}_0\,\deleted{\unsure{(}}S\deleted{\unsure{)}}\longrightarrow\underline{\mathrm{Cofib}}_0\,S
\]
% L13-15
induisent des équivalences sur les localisées (chose que j'ai été incapable de prouver tantôt)$^{*}$. \margin{$^{*}$ Les \unsure{quasi-inv.} étant bel et bien donnés par $\Phi_0,\Psi_0$ -- car par quoi d'autre ?}
% slanted
Par contre, il est \inserted{pire que} déraisonnable de s'attendre \inserted{ici} qu'il en soit de $=$ pour les foncteurs
% L16-17
\[
 \underline{\mathrm{Fibsc}}_0\,S\longrightarrow\underline{\mathrm{Fibsc}}\,S,\quad
 \underline{\mathrm{Cofibsc}}_0\,S\unsure{\mathord{-}|\longrightarrow}\underline{\mathrm{Cofibsc}}\,S
\]
% L18-21
(ça devrait donner par localisation des foncteurs pleinement fidèles d'inclusion, des catégories des objets loc. const. dans $\mathrm{HOT}_W(S)$ resp. dans $(\unsure{\mathrm{HOt}|\mathrm{HOT}}_W(S^{\mathrm o}))^{\mathrm o}$. Et de même pour
% archive sheet number: 88
\scanpage{90}\mspage{34}
% L1-2
pour les foncteurs induits sur les loc. par les
\[
 \underline{\mathrm{Fibsc}}\,S\longrightarrow\underline{\mathrm{Fib}}\,S,\quad
 \underline{\mathrm{Cofibsc}}\,S\longrightarrow\underline{\mathrm{Cofib}}\,S.
\]
% L3-7
Comme les localisées des catégories buts des ces deux foncteurs ont tout l'air de s'identifier aux sous-catégories précédentes \deleted{\unsure{d}}d'objets loc constants, il semblerait donc que \textcircled{\scriptsize 2\textsuperscript{o}} les foncteurs précédents définissent
% L8-11
des rétractions de la catégorie $\mathrm{HOT}_W(S)$ resp. $\mathrm{HOT}_W(S^{\unsure{\mathrm o|\circ}})^{\mathrm o}$ sur la sous-catégorie des objets loc. constants sur $S$, les foncteurs localisés des foncteurs d'inclusion
% L12
\[
 \underline{\mathrm{Fibsc}}_0\,S\longrightarrow\underline{\mathrm{Fibsc}}\,S,\quad
 \underline{\mathrm{Cofibsc}}_0\,S\longrightarrow\underline{\mathrm{Cofibsc}}\,S,
\]
% L13-14
étant pleinement fidèles et ayant comme images essentielles les sous-catégories en question.

% L15-20
Cette rétraction m'apparaît comme une chose tout à fait mystérieuse. \deleted{Pourtant, il semblerait que pour \unsure{le} \deleted{\unsure{dérivé}} dérivateur $\mathbb D$ (pas seulement HOT) qui provient \unsure{d'une|d'un} $(M,W)$ \deleted{avec $M$} provenant d'un \uline{triple de Quillen},} C'est lié au fait que pour le dérivateur très spécial $\mathbb D=\mathrm{HOT}$,
% L21-22
quand on a un $\xi\in\mathbb D(S)$, on en déduit \deleted{\unsure{une}} (à $\unsure{\mathrm{HOT}_{W/S}|\mathrm{HOT}_{/S}}$-équivalence près) un \deleted{$\unsure{\xi}$} \deleted{\unsure{type de}} \inserted{diagramme $X$ au dessus de $S$,} \deleted{\unsure{en quelque}}
% archive sheet number: 89
\scanpage{91}\mspage{35}
% L1-2
qu'on peut regarder \unsure{donc|alors} \inserted{(via la construction $C$)} comme définissant un type d'homotopie relatif \uline{loc constant} sur $S$.

% L3-7
Donc pour élucider la situation i.e. les relations entre les 13 catégories localisées provenant des diagrammes de la page 8, il reste à regarder si les assertions hypothétiques que je viens de faire sont bien valables, savoir

% L8-9
\textcircled{\scriptsize 1\textsuperscript{o}} Le foncteur déduit de
\[
 \underline{\mathrm{Fibsc}}_0\,\deleted{\unsure{E}}S\longrightarrow\underline{\mathrm{Fib}}_0\,S
\]
% L10-15
par passage aux localisés est une équivalence de catégories (et on peut dans cette question remplacer $\underline{\mathrm{Fib}}_0\,S$ par n'importe quelle \inserted{autre} sous catégorie pleine de $\underline{\mathrm{Cat}}_{/S}$ la contenant -- cela revient au même), et

% L16-17
\textcircled{\scriptsize 2\textsuperscript{o}} Le foncteur déduit de
\[
 \underline{\mathrm{Fibsc}}_0\,\deleted{\unsure{E}}S\longrightarrow\underline{\mathrm{Fibsc}}\,S
\]
% L18-22
par localisation est pleinement fidèle. (Que son image \deleted{essentielle} soit exactement formée des objets loc. \unsure{constants} \unsure{du premier|des premiers} \unsure{localisé|localisés}, interprété comme $\mathrm{HOT}_W(S)$, est \deleted{\unsure{doute}} logique.)\unsure{.}

% L23
Il n'y a plus lieu de se soucier du côté
% archive sheet number: 90
\scanpage{92}\mspage{36}
% L1-3
cofibrés, puisque tout ce qu'on \unsure{pourra|pourrait} en dire revient à ce que je viens d'expliciter, pour $X^{\unsure{\mathrm o|\circ}}$ au lieu de $X$.

% L4-6
Dans la direction de 1\textsuperscript{o}, les arguments précédents à coups de $\Phi_0,p$ nous donnent déjà un renseignement utile.
% L7-11
Mais il sera plus commode à présent de laisser tomber $\underline{\mathrm{Fib}}_0\,S$ ou plutôt, de le remplacer par $\underline{\mathrm{Fib}}\,S$ (ce qui est licite, on l'a dit dans 1\textsuperscript{o}), de sorte qu'on a maintenant un diagramme commutatif (traits pleins)
% L12-14
\diagram{\begin{tikzcd}[ampersand replacement=\&,column sep=small]
 \underline{\mathrm{Fibsc}}\,S\arrow[dr,"j"']\&\&\underline{\mathrm{Fibsc}}_0\,S\arrow[ll,hook,"i"']\arrow[dl,"k"']\\
 \&\underline{\mathrm{Fib}}\,S\arrow[ur,dashed,bend right=38,"{\Phi_0=\unsure{\varphi|\psi}}"']\&
\end{tikzcd}}
,
% L15-18
où $i$ est une inclusion (strictement) pleinement fidèle, $j$ (donc aussi $k$) fidèle. Enfin, on a un foncteur canonique $\varphi(=\Phi_0)$ \deleted{\unsure{tel que}} (en pointillés). Les localis\deleted{\unsure{ations}}eurs sur la ligne du haut sont les localiseurs fins, (pour $\underline{\mathrm{Fibsc}}_0\,S$
% L19-23
c'est pareil que la $W$-équivalence grossière), sur $\underline{\mathrm{Fib}}\,S$ par contre on prend la $W$-équivalence grossière. (On sait que le localisé est le $=$, à équivalence près, que si on se restreint à la sous-catégorie pleine $\underline{\mathrm{Fib}}_0\,S$, ce qui est assez remarquable.)
% archive sheet number: 91
\scanpage{93}\mspage{37}
% L1-2
\deleted{\unsure{D'aut}} Tous les \inserted{quatre} foncteurs $i,j,k,\Phi_0$ sont compatibles aux localiseurs, donc passent aux localisés
% L3-5
\diagram{\begin{tikzcd}[ampersand replacement=\&,column sep=small]
 \mathrm{HOT}_W(S)\arrow[dr,"{\bar j}"']\&\& \mathrm{HOT}^{?}_W(S)\arrow[ll,"{\bar i}"']\arrow[dl,"{\bar k}"']\\
 \&\mathrm{HOT}^{\mathrm{lc}}_W(S)\arrow[ur,bend right=32,"{\bar{\unsure{\varphi|\psi}}}"']\&
\end{tikzcd}}
,
% L6-9
où l'exposant $?$ dans la notation signifie que je ne sais interpréter cette catégorie, sauf si je prouve que $\bar k$ est une équivalence\unsure{;|,} et l'exposant $\mathrm{lc}$ dans la catégorie du bas
% L10-14
doit suggérer \deleted{\unsure{ma présomption}} \inserted{l'intuition} qu'il s'agit \deleted{\unsure{d'} la \unsure{sous-c}} en tous cas d'une sorte de ``types d'homotopie relatifs'' (dans un sens \unsure{modifié}) \uline{localement constants} sur $S$. Le problème est de comprendre ce diagramme.

% L15-21
La conjecture naturelle est

1\textsuperscript{o}) $\bar k$ est une équivalence. (Ce qui impliquera que $\bar\varphi$ est l'équivalence quasi-inverse\unsure{)}, que \deleted{2\textsuperscript{o})} $\bar k$, \inserted{$\bar\varphi$} permettent \unsure{essentiellement} d'identifier $\mathrm{HOT}^{?}_W$ et $\mathrm{HOT}^{\mathrm{lc}}_W$, de sorte que $\bar i$ et $\bar j$ deviennent des foncteurs entre $\mathrm{HOT}_W$ et $\mathrm{HOT}^{\unsure{\mathrm{loc}|\mathrm{lc}}}_W$.)

% L22-24
2\textsuperscript{o}) $\bar i$ est pleinement fidèle (NB l'image essentielle est bien comprise de \unsure{toute façons|toute façon}.)
% archive sheet number: 92
\scanpage{94}\mspage{38}
% L1-3
En direction de 1\textsuperscript{o}), je sais pour l'instant seulement que l'on a
\[
 \bar k\,\bar{\unsure{\varphi|\psi}}\simeq\mathrm{id}_{\mathrm{HOT}^{\mathrm{lc}}_W(S)}
\]
% L4-8
donc $\bar k$ est ess. surjectif (sur les \deleted{\unsure{f}} objets), et induit des \uline{surjections} sur les $\mathrm{Hom}(-,-)$, tandis que $\bar\varphi\deleted{\unsure{\Phi}}$ est injectif sur les classes d'isomorphie d'objets, et sur les $\mathrm{Hom}(-,-)$, \deleted{\unsure{cela}} (i.e. \uline{fidèle}).

% L9-13
Si les relations entre $\mathrm{HOT}^{?}_W$ et $\mathrm{HOT}^{\mathrm{lc}}_W$ s'avèrent \unsure{indétables|indésirables|indébrouillables} (p.ex. $\bar\varphi$ pas $=$ pl. fid i.e. $\bar k$ pas même fidèle), alors la catégorie $\mathrm{HOT}^{?}_W$ mérite d'être larguée dans la corbeille, et on ne gardera que
% L14-16
\diagram{\begin{tikzcd}[ampersand replacement=\&,row sep=large]
 \mathrm{HOT}_W(S)\arrow[d,shift right=2,"\bar j"']\\
 \mathrm{HOT}^{\mathrm{lc}}_W(S)\arrow[u,shift right=2,"{\bar{\unsure{\psi|\Psi}}'}"'],
\end{tikzcd}}
% L17-18
où on définit $\bar\psi'$ comme
\[
 \bar\psi'=\bar i\,\bar\varphi=\overline{i\deleted{\unsure{\Phi}}\varphi}.
\]
% L19-20
Et la question \inserted{(plus modeste)} devient à présent :

3\textsuperscript{o}) $\bar\psi'$ est-il pleinement fidèle ?
% L21-24
(L'image essentielle est \deleted{toujours la} \inserted{la $=$} que celle de $\overline{\underline{\mathrm{Fib}}_0}$ \deleted{par le \unsure{côté}} $(\approx\overline{\underline{\mathrm{Fib}}})$, elle est formée \deleted{des} sauf erreur des objets loc constants -- à vérifier au besoin\dots)
% archive sheet number: 93
\scanpage{95}\mspage{39}
% L1-3
Il est difficile d'imaginer que $\bar k$, ou ce qui revient au $=$ $\bar{\unsure{\varphi|\psi}}$, puisse être une équivalence, sans qu'il existe au moins une ``$W_{\unsure{f|S}}$-homotopie''
% L4
\diagram{\begin{tikzcd}[ampersand replacement=\&,column sep=0.43em]
 \deleted{\unsure{\text{entre}}}\ \deleted{\unsure{\varphi k}}\quad
 \mathrm{id}_{\underline{\mathrm{Fibsc}}_{\deleted{\unsure{?}}0}(\deleted{\unsure{E}}S)}
 \arrow[r]\arrow[rrrr,phantom,"h"]\&{}\&{}\arrow[l]\arrow[r]\&{}\&{}\arrow[l]\arrow[r]\&{}\&{}\arrow[l]\arrow[r]\&\deleted{\unsure{\bar{\phantom{k}}}}\varphi k
\end{tikzcd}}
% L5-7
dans la catégorie des endofoncteurs de $\underline{\mathrm{Fibsc}}_0\,S$. \deleted{Donc il faut, pour $X$ objet donné dans cette catégorie,}

% L7-12
À vrai dire, je me rends compte maintenant que la description de $\bar\varphi$ à partir de $\varphi:X\longmapsto\Phi_0 X=\Phi C(X)$ est inutilement compliquée, car on avait fait le détour par $C(X)$ uniquement pour être sûr de tomber dans $\underline{\mathrm{Fibsc}}_0\,S$, pas seulement dans $\underline{\mathrm{Fibsc}}$.
% L13-17
Mais si on laisse tomber $\underline{\mathrm{Fibsc}}_0\,S$, on s'en fout, et le foncteur $\unsure{\varphi'|\varphi}:X\mapsto\Phi(X)$ est beaucoup plus naturel. Et \deleted{\unsure{si}} si on ne laisse pas tomber, il n'y a pas d'inconvénient de remplacer la source $\underline{\mathrm{Fib}}\,S$
% L18-22
par $\underline{\mathrm{Fib}}_0\,S$ (lequel a $=$ localisé, à équi. près), et sur $\underline{\mathrm{Fib}}_0\,S$ à nouveau il n'y a pas lieu de s'encombrer de $CX$, pour le plaisir de lui appliquer $\Phi$, car $\Phi(X)$ lui-même est déjà dans $\underline{\mathrm{Fibsc}}_0\,S$. Donc, la question est celle-ci : Si $X$ varie
% archive sheet number: 94
\scanpage{96}\mspage{40}
% L1-2
dans $\underline{\mathrm{Fibsc}}_{\underset{\uparrow\uparrow\nwarrow}{0}}\,S$, y-a-t-il une homotopie dans $\underline{\mathrm{Fibsc}}_{\underset{\uparrow\uparrow\uparrow}{0}}\,S$, ``naturelle'' i.e. fonctorielle en $\unsure{x|X}$,
% L3-5
\diagram{\begin{tikzcd}[ampersand replacement=\&,column sep=small]
 X\arrow[r]\&{}\arrow[r,"{\deleted{\unsure{\Phi}}\inserted{c}}",no head]\&
 \Phi(X)=\begin{gathered}\text{Catégorie \deleted{\unsure{co}}\unsure{fibrée|cofibrée}}\\
 \text{scindée de \unsure{fibre|fibres}}\\ \text{les }{}_{s\backslash}X\end{gathered}
\end{tikzcd}}
% L6-7
Mais un chemin qui soit formé de flèches
\[
 X=X_0\mathbin{-}X_1\mathbin{-}X_2\ldots\mathbin{-}X_n=\Phi(X)
\]
% L8-12
qui soient toutes \uline{cartésiennes} (et \unsure{de} plus, bien \unsure{sur|sûr}, dans $W$ !). Cela signifie donc qu'on \deleted{a des \unsure{c}} aurait des \deleted{\unsure{esp}} cat. fibrées scindées sur $S$, $X_1,\ldots,X_{n-1}$, dépendant \uline{fonctoriellement} de $X$, qui relient
% L13-14
\[
 X=X_0\qquad \text{à}\qquad \Phi(X)=X_n
\]
i.e.
\[
 X_s\qquad\text{à}\qquad {}_{s\backslash}X_{\deleted{\unsure{/s}}}\qquad ???
\]
% L15
Je ne vois rien à l'horizon
% archive sheet number: 95
\scanpage{97}\mspage{41}
% archive sheet number: 96
% L1
\unsure{15.11.|15,11.}

% L2-5
Je crois que j'ai enfin compris, cette nuit, la situation exacte. Il y a lieu, dans $\mathrm{Cat}_{/S}$ et dans les autres catégories des diagrammes p.~8 (s'envoyant toutes dans $\mathrm{Cat}_{/S}$)\unsure{,} d'envisager \deleted{trois} \inserted{quatre} sortes de ``$W$-équivalence''.

% L6-8
a) La $W$-\uline{équivalence} ``\uline{grossière}'' :
\diagram{\begin{tikzcd}[ampersand replacement=\&,column sep=tiny,row sep=tiny]
 f:X\arrow[rr]\arrow[dr,no head]\&\&Y\arrow[dl,no head]\\
 \&S\&
\end{tikzcd}}
\inserted{dans $\mathrm{Cat}_{/S}$} est une $W$-équivalence grossière si en \deleted{\unsure{tant}} que flèche de $\mathrm{Cat}$, c'est une $W$-équivalence.
\margin{Stable par chgt de base à la fois lisse et propre, plus \unsure{génér.} si $S'\longrightarrow S$ est une Hot-fibration}
% slanted

% L9-10
b) La $W$-\uline{équivalence} ``\uline{par fibres}\inserted{(sur $S$)}'' : $\forall s\in S$, $f_s:X_s\longrightarrow Y_s$ est une $W$-équivalence. \deleted{(Signifie que c'est une $W$-équiv, et le reste par \unsure{tt chgt}}
\margin{Stable par tt chgt de base $S'\longrightarrow S$}
% slanted

% L11-12
c) La \uline{$W$-équivalence} ``\deleted{\unsure{co}}\uline{locale}'' \uline{sur} $S$'' : $\forall s\in S$, $f_{/s}:X_{/s}\longrightarrow Y_{/s}$ est une $W$-équivalence. \inserted{Signifie que c'est une $W$-équiv., et le reste par \deleted{tout} \unsure{chgt} de base $S'\longrightarrow S$ qui est un isom. local.}
\margin{devrait être stable par chgt de base \uline{lisse} $S'\longrightarrow S$}
% slanted

% L13-14
\deleted{\unsure{d)}}\unsure{c')|d)} La $W$-\uline{équivalence} ``\uline{colocale sur} $S$'' : $\forall s\in S$, $ {}_{s\backslash}f:{}_{s\backslash}X\longrightarrow{}_{s\backslash}Y$ est une $W$-équivalence. \inserted{Signifie que c'est une $W$-équivalence et le reste par tt chgt. de base $S'\longrightarrow S$ qui est un co-isom local.}
\margin{devrait être stable \unsure{pour|par} chgt. de base \uline{propre} sur $S$}
% slanted

% L15-19
Je vais résumer dans un énoncé récapitulatif les relations entre ces quatre notions de $W$-équivalences, en pensant au cas $W=W_{\unsure{\infty|0}}$. Une fois la chose faite, je me préoccuperai de dégager les conditions minimales sur \inserted{un} $W$ \inserted{général} qui assurent que c'est vrai pour $W$.

% L20
\uline{Proposition 1} \unsure{.} Soit
\diagram{\begin{tikzcd}[ampersand replacement=\&,column sep=tiny,row sep=tiny]
 f:X\arrow[rr]\arrow[dr,no head]\&\&Y\arrow[dl,no head]\\
 \&S\&
\end{tikzcd}}
dans $\mathrm{Cat}_{/S}$.

% L21-23
1) La $W$-équivalence locale sur $S$, et la $W$-équivalence colocale sur $S$, impliquent la $W$-équivalence grossière :
\diagram{\begin{tikzcd}[ampersand replacement=\&,row sep=tiny,column sep=tiny]
 \deleted{\unsure{a}}c\arrow[dr,Rightarrow]\&\&c'\deleted{\unsure{d}}\arrow[dl,Rightarrow]\\
 \&a\&
\end{tikzcd}}

% L24-27
2) Si $X,Y$ sont propres sur $S$, ou $X,Y$ lisses sur $S$, ou $X,Y$ des Hot-fibrés sur $S$, alors la $W$-équivalence par fibres implique que c'est une $W$-équivalence \uline{universelle} (\inserted{sur $S$}) (\deleted{stable} $W$-équivalence, et le reste par \uline{tt} chg. de base $S'\longrightarrow S$) donc dans ces \unsure{trois} cas, on a
\margin{Il suffit $X,Y$ $W$-propres sur $S$, ou $W$-lisses sur $S$, ou $W$-fibrés sur $S$,}
% slanted
\diagram{\begin{tikzcd}[ampersand replacement=\&,column sep=tiny,row sep=tiny]
 \& b\arrow[dl,Rightarrow]\arrow[dr,Rightarrow]\&\\
 c\arrow[dr,Rightarrow]\&\&c'\arrow[dl,Rightarrow]\\
 \&a\&
\end{tikzcd}}
% L28-31
Mais on peut préciser \unsure{chacun} \deleted{\unsure{cas}} de ces trois cas :

\unsure{--}(2 fib) Si $X,Y$ sont Hot-fibrés sur $S$, alors la $W$-équivalence \inserted{gross.} de $f$ implique que c'est une $W$-équiv. par fibres, donc les quatre notions de $W$-équivalence sont équivalentes pour $f$.
\margin{$a\Rightarrow b$ seulement moyennant L~5, voir ci-dessous p.~43}
% slanted
\diagram{\begin{tikzcd}[ampersand replacement=\&,column sep=tiny]
 \unsure{a}\&\unsure{\Longleftarrow|\Longleftrightarrow}\&c\&\Longleftrightarrow\&c'\&\Longleftarrow\&b
\end{tikzcd}}
\scanpage{98}\mspage{42}
% archive sheet number: 97
% L1-4
(2 lisse) Si $X,Y$ sont lisses sur $S$, alors la $W$-équivalence par fibres sur $S$ équivaut à la $W$-équivalence colocale sur $S$, \inserted{et implique la $W$-équiv. \uline{universelle} sur $S$,} donc dans ce cas
\margin{il suffit $=$ $X,Y$ $W$-lisses sur $S$}
% slanted
\diagram{\begin{tikzcd}[ampersand replacement=\&,row sep=tiny,column sep=small]
 b\arrow[r,Leftrightarrow]\arrow[d,Rightarrow]\&c'\\
 c\arrow[d,Rightarrow]\&\\
 a\&
\end{tikzcd}}

% L5-7
(2 \unsure{prop|propr}) Si $X,Y$ sont propres sur $S$, alors la $W$-équiv. par fibres sur $S$ équivaut à la $W$-équivalence locale sur $S$ :
\margin{Il suffit $=$ $X,Y$ $W$-propres sur $S$}
% slanted
\diagram{\begin{tikzcd}[ampersand replacement=\&,row sep=tiny,column sep=small]
 b\arrow[r,Leftrightarrow]\arrow[d,Rightarrow]\&c\\
 c'\arrow[d,Rightarrow]\&\\
 a\&
\end{tikzcd}}

% L8-15
\uline{Dém} 1) n'est autre que L~4 et \deleted{\unsure{L 4 et L 4 bis}} son dual L~4\unsure{$'$|$'{}^{*}$}, pour un localiseur fond. $W$, et on a vu que si on suppose L~1, L~3, alors L~4 équivaut à L~4$'$, et dans la notion de loc. fond. on supposera tjrs que c'est vérifié. Également, nous supposons que L~3 est vraie $=$ sous la forme plus forte L~3 \unsure{quater} (p.~1) : $W_\infty\subset W$, i.e., toute Hot-équivalence \deleted{\unsure{est}}est une $W$-équivalence.
\margin{Mais il y a confusion L~4 est plus faible, cf p.~\unsure{4 §|44|45} l'axiome L~4 bis}
% slanted

% L16-18
(2 \deleted{\unsure{liss}}\unsure{lisse}) Si $X,Y$ lisses sur $S$, alors les $X_s\longrightarrow{}_{s\backslash}X$ $Y_s\longrightarrow{}_{s\backslash}Y$ sont des $W_\infty$-équiv\unsure{,|.} \inserted{donc des $W$-équiv. (il suffirait $=$ que $X,Y$ \deleted{\unsure{soient}} \inserted{\unsure{soient}} $W$-lisses/$S$)} donc le diagramme comm.
\diagram{\begin{tikzcd}[ampersand replacement=\&,row sep=small,column sep=large]
 X_s\arrow[r,"f_s",no head]\arrow[d,"i_s^X"']\&Y_s\arrow[d,"i_s^Y"]\\
 {}_{s\backslash}X\arrow[r,"{}_{s\backslash}f"]\&{}_{s\backslash}Y
\end{tikzcd}}
% L19-24
nous montre que $f_s$ est une $W$-équiv. ssi $ {}_{s\backslash}f$ l'est, d'où $b\Longleftrightarrow c'$. Pour voir que la $W$-équivalence par fibres est \inserted{$=$} une $W$-équivalence \uline{universelle} sur $S$, \deleted{\illegible{1}} on est ramené à montrer que c'est une $W$-équiv. grossière (puisque \deleted{\unsure{universelle}} la lissité stable par chgt. de base sur $S$), \deleted{\unsure{en fait}} \inserted{\unsure{et que|et par cons. que}} \deleted{il devrait suffire $=$ que $X,Y$ soient $W$-lisses sur $S$, \unsure{ou} $W$ \unsure{au} des lisses.} puisque $c'\Rightarrow a$ par 1).

% L25-26
Le cas (2 propre) est dual.

(2 \deleted{\unsure{fibr}}\unsure{fib}) Si $X,Y$ sont des Hot-fibrés sur $S$, alors la
